\section{Experimentos}

\noindent La experimentación tuvo como objetivo evaluar el comportamiento del algoritmo propuesto para el MS-CFLP-CI en términos de calidad de solución, factibilidad y tiempo de ejecución. Para ello, se utilizó una implementación en C++ basada en Hill Climbing con estrategia Best Improvement y reinicios. El algoritmo trabaja sobre una representación mediante matriz de envíos \(x_{ji}\), lo que permite modificar las asignaciones entre clientes y bodegas manteniendo consistencia con la variante multi-source del problema.\\

\noindent Las ejecuciones se realizaron en un computador Acer Predator Triton 300 SE, modelo Predator PT314-51s, con procesador 11th Gen Intel(R) Core(TM) i7-11375H @ 3.30GHz, \(16\) GB de memoria RAM y sistema operativo de \(64\) bits. Para mantener una comparación consistente, todas las instancias fueron ejecutadas utilizando el mismo código fuente, compilador y configuración experimental. Como indicadores de evaluación se consideraron el valor final de la función objetivo, la factibilidad de la solución obtenida, el tiempo de ejecución y la mejora alcanzada respecto de una solución inicial greedy.\\

\noindent El procedimiento experimental consideró tres configuraciones principales. Primero, se evaluó una construcción greedy determinística, equivalente a ejecutar el algoritmo con \(\epsilon = 0\), un solo reinicio y sin iteraciones de búsqueda local. Esta configuración se utilizó como línea base. Luego, se ejecutó Hill Climbing Best Improvement desde una solución inicial determinística, con \(\epsilon = 0\), un reinicio y un número fijo de iteraciones. Finalmente, se evaluó la versión estocástica con reinicios, donde cada reinicio genera una solución inicial mediante una regla \(\epsilon\)-greedy.\\

\noindent Para todas las configuraciones principales se utilizaron cinco semillas: \(101\), \(202\), \(303\), \(404\) y \(505\). La versión con \(\epsilon = 0\) es determinística, ya que bajo las mismas condiciones construye la misma solución inicial y aplica el mismo criterio de mejora. En cambio, la versión con \(\epsilon > 0\) y reinicios es estocástica, debido a que distintas semillas pueden generar soluciones iniciales diferentes. Esta diferencia permite analizar el efecto de la diversificación en la calidad de las soluciones obtenidas.\\

\noindent Para estudiar el efecto del parámetro \(\epsilon\), se utilizó la instancia wlp01.in, compuesta por \(50\) bodegas y \(115\) clientes. Esta instancia fue seleccionada por tener un tamaño intermedio pequeño, lo que permite ejecutar varias semillas en tiempos razonables. Se probaron los valores \(\epsilon \in \{0, 0.10, 0.20, 0.40\}\), con \(5\) reinicios, \(30\) iteraciones máximas por reinicio y las cinco semillas definidas. A partir de estos resultados se seleccionó \(\epsilon = 0.20\), ya que entregó el menor costo promedio.\\

\noindent Posteriormente, se estudió el desempeño del algoritmo en instancias de distinto tamaño. Se utilizaron las instancias toy.in, wlp01.in, wlp02.in, wlp03.in, wlp04.in, wlp05.in, wlp06.in y wlp07.in, cubriendo más de cinco casos de prueba. Para toy.in, wlp01.in, wlp02.in y wlp03.in se utilizó la configuración \(\epsilon = 0.20\), \(5\) reinicios y \(30\) iteraciones por reinicio. Para las instancias más grandes, desde wlp04.in hasta wlp07.in, se redujo la configuración a \(3\) reinicios y \(20\) iteraciones por reinicio, con el objetivo de mantener tiempos de ejecución razonables debido al crecimiento del vecindario.\\

\noindent Además, se realizó un análisis de convergencia sobre la instancia wlp03.in variando el número máximo de iteraciones por reinicio en \(\{0, 5, 10, 20, 30\}\). Este experimento permitió observar cómo cambia la calidad de la solución respecto del avance de la búsqueda y del tiempo de ejecución.\\

\noindent El criterio de término del algoritmo estuvo dado por el número máximo de reinicios y el número máximo de iteraciones por reinicio. Además, se consideró un límite máximo de tiempo por ejecución para evitar corridas excesivamente largas. En cada ejecución se registró la mejor solución encontrada, su costo total, su factibilidad y el tiempo requerido. De esta forma, el diseño experimental permite comparar la línea base greedy con la búsqueda local determinística, la versión estocástica con reinicios y el comportamiento de convergencia del método.